% SOURCE AUTHORITY:
% expose_VIII_body.tex lines 2316--2361; scan indices 315--317;
% source folios 304--306.
% Direct scan comparison: physical PDF pages 316--318 in overlapping
% 1100-dpi bands. The kernel formulas and the symmetry diagram were checked
% directly; all nodes, arrows, label sides, and the closing comma are clear.
% Transparent source dispositions:
% - typewritten letter-o special-fibre subscripts are normalized to zero;
% - the biextension denoted once by E rather than E_eta in 7.3.2 is rendered
%   consistently as the stated generic-fibre biextension E_eta;
% - "biextension symetrique" in the construction of {}^sE is rendered
%   "transposed biextension" to distinguish that operation from the later
%   property E_eta \simeq {}^sE_eta of being symmetric.
% Hard stop before 7.3.4, source line 2367 / scan index 317.

\medskip
\noindent 7.3.2.  It follows from the compatibility (7.3.1.1) that, if
$E_\eta$ is an extension of $P_\eta$ by $\underline G_{m\eta}$, there is a
largest group subscheme $V$ of $P$ such that the restriction of $E_\eta$ to
$V_\eta$ extends to an extension of $V$ by $\underline G_{mS}$.  It is the
open subgroup $V$ of $P$ containing $P_\eta$ and defined by the condition
\begin{equation}\tag{7.3.2.1}
 V_0/V_0^{\circ}=\ker d(E_\eta).
\end{equation}
If $u:P'\to P$ is then a morphism of smooth group schemes over $S$, the
extension $E'_\eta$, obtained by pulling $E_\eta$ back along $u_\eta$,
extends to an extension of $P'$ by $\underline G_{mS}$ if and only if
\begin{equation}\tag{7.3.2.2}
 u(P')\subset V.
\end{equation}

In the case of a biextension $E_\eta$ of $(P_\eta,Q_\eta)$ by
$\underline G_{m\eta}$, there is likewise a largest group subscheme $V$ of
$P$ such that the restriction of $E_\eta$ to $(V_\eta,Q_\eta)$ extends to
a biextension of $(V,Q)$ by $\underline G_{mS}$.  It is the open subgroup
of $P$ containing the generic fibre $P_\eta$ and characterized by the fact
that $V_0/V_0^{\circ}$ is the left annihilator in $\Phi$ of the pairing
(7.1.2), that is, the kernel of the corresponding homomorphism from $\Phi$
to $\underline{\mathrm{Hom}}(\Psi,(\mathbb Q/\mathbb Z)_k)$:
\begin{equation}\tag{7.3.2.3}
 V_0/V_0^{\circ}
 =\ker\!\left(
   \Phi\xrightarrow{\ d(E_\eta)\ }
   \underline{\mathrm{Hom}}(\Psi,(\mathbb Q/\mathbb Z)_k)
 \right).
\end{equation}
If $u:P'\to P$ is a morphism of smooth group schemes over $S$, the
biextension of $(P'_\eta,Q_\eta)$ by $\underline G_{m\eta}$ obtained by
pulling $E_\eta$ back along $u_\eta$ extends to a biextension of $(P',Q)$
by $\underline G_{mS}$ if and only if (7.3.2.2) holds.  Symmetric results
hold upon interchanging the roles of $P$ and $Q$.

Of course, in general there is no largest \underline{pair} of subgroups
$P'\to P$, $Q'\to Q$ for which the restriction $E'_\eta$ of $E_\eta$ to a
biextension of $(P'_\eta,Q'_\eta)$ by $\underline G_{m\eta}$ extends to a
biextension of $(P',Q')$ by $\underline G_{mS}$.  One can only say that,
for a given pair $(P',Q')$, or more generally for a given pair of
morphisms $P'\to P$, $Q'\to Q$ of smooth group schemes, the desired
extension exists if and only if the restriction of the form $d(E_\eta)$
to
\[
 P'_0/{P'_0}^{\circ}\otimes Q'_0/{Q'_0}^{\circ}
\]
is zero.

\medskip
\noindent 7.3.3.  \underline{Symmetry}.  Let $E_\eta$ be a biextension of
$(P_\eta,Q_\eta)$ by $\underline G_{m\eta}$, and consider the transposed
biextension ${}^{s}E_\eta$ (VII 2.7), which is a biextension of
$(Q_\eta,P_\eta)$ by $\underline G_{m\eta}$.  The following diagram is
commutative:
\begin{equation}\tag{7.3.3.1}
\begin{tikzcd}[row sep=large,column sep=huge]
 \Phi\otimes\Psi
   \arrow[dd,"\mathrm{sym}"']
   \arrow[rd,"d(E_\eta)"] & \\
 & (\mathbb Q/\mathbb Z)_k \\
 \Psi\otimes\Phi
   \arrow[ru,"d({}^{s}E_\eta)"'] &
\end{tikzcd}
\qquad ,
\end{equation}
where the vertical arrow is the canonical symmetry.  (Compare this
statement with the \underline{antisymmetry} statement of 2.2.10.)  It
follows in particular that, if $P=Q$ and $E_\eta$ is a
\underline{symmetric} biextension of $(P_\eta,Q_\eta)$ by
$\underline G_{m\eta}$, that is, if $E_\eta$ is isomorphic to
${}^{s}E_\eta$, then the bilinear form $d(E_\eta)$ on $\Phi\times\Phi$
with values in $(\mathbb Q/\mathbb Z)_k$ is \underline{symmetric} (and not
alternating, as in 2.2.12!).  If it is zero, that is, if $E_\eta$ extends
to a biextension $E$ of $(P,P)$ by $\underline G_{mS}$, the latter is
necessarily symmetric as well: by 7.1(a), every isomorphism
${}^{s}E_\eta\to E_\eta$ extends uniquely to an isomorphism
${}^{s}E\to E$.
